\documentclass[10pt,a4paper]{amsart}
\usepackage{amsmath,amssymb,amsthm,graphicx}
% Retained minimal preamble for the Stacks Project chapter more-morphisms.tex (More on Morphisms).
% Verbatim excerpt of the chapter preamble: Xy-pic setup (native preamble lines 14-18) and the native
% theorem declarations with the mathematical macros (native preamble lines 154-217). Copied mechanically;
% nothing added, removed or redefined. The wrapper class and the amsmath/amssymb/amsthm/graphicx base
% packages are supplied by the assembler, so the class/externaldocument part of the native preamble is omitted.
% For commutative diagrams we use Xy-pic
\usepackage[all]{xy}

% We use 2cell for 2-commutative diagrams.
\xyoption{2cell}
% Theorem environments.
%
\theoremstyle{plain}
\newtheorem{theorem}[subsection]{Theorem}
\newtheorem{proposition}[subsection]{Proposition}
\newtheorem{lemma}[subsection]{Lemma}

\theoremstyle{definition}
\newtheorem{definition}[subsection]{Definition}
\newtheorem{example}[subsection]{Example}
\newtheorem{exercise}[subsection]{Exercise}
\newtheorem{situation}[subsection]{Situation}

\theoremstyle{remark}
\newtheorem{remark}[subsection]{Remark}
\newtheorem{remarks}[subsection]{Remarks}

\numberwithin{equation}{subsection}

% Macros
%
\def\lim{\mathop{\mathrm{lim}}\nolimits}
\def\colim{\mathop{\mathrm{colim}}\nolimits}
\def\Spec{\mathop{\mathrm{Spec}}}
\def\Hom{\mathop{\mathrm{Hom}}\nolimits}
\def\Ext{\mathop{\mathrm{Ext}}\nolimits}
\def\SheafHom{\mathop{\mathcal{H}\!\mathit{om}}\nolimits}
\def\SheafExt{\mathop{\mathcal{E}\!\mathit{xt}}\nolimits}
\def\Sch{\mathit{Sch}}
\def\Mor{\mathop{\mathrm{Mor}}\nolimits}
\def\Ob{\mathop{\mathrm{Ob}}\nolimits}
\def\Sh{\mathop{\mathit{Sh}}\nolimits}
\def\NL{\mathop{N\!L}\nolimits}
\def\CH{\mathop{\mathrm{CH}}\nolimits}
\def\proetale{{pro\text{-}\acute{e}tale}}
\def\etale{{\acute{e}tale}}
\def\QCoh{\mathit{QCoh}}
\def\Ker{\mathop{\mathrm{Ker}}}
\def\Im{\mathop{\mathrm{Im}}}
\def\Coker{\mathop{\mathrm{Coker}}}
\def\Coim{\mathop{\mathrm{Coim}}}

% Boxtimes
%
\DeclareMathSymbol{\boxtimes}{\mathbin}{AMSa}{"02}

%
% Macros for moduli stacks/spaces
%
\def\QCohstack{\mathcal{QC}\!\mathit{oh}}
\def\Cohstack{\mathcal{C}\!\mathit{oh}}
\def\Spacesstack{\mathcal{S}\!\mathit{paces}}
\def\Quotfunctor{\mathrm{Quot}}
\def\Hilbfunctor{\mathrm{Hilb}}
\def\Curvesstack{\mathcal{C}\!\mathit{urves}}
\def\Polarizedstack{\mathcal{P}\!\mathit{olarized}}
\def\Complexesstack{\mathcal{C}\!\mathit{omplexes}}
% \Pic is the operator that assigns to X its picard group, usage \Pic(X)
% \Picardstack_{X/B} denotes the Picard stack of X over B
% \Picardfunctor_{X/B} denotes the Picard functor of X over B
\def\Pic{\mathop{\mathrm{Pic}}\nolimits}
\def\Picardstack{\mathcal{P}\!\mathit{ic}}
\def\Picardfunctor{\mathrm{Pic}}
\def\Deformationcategory{\mathcal{D}\!\mathit{ef}}
\newenvironment{reference}{}{}
\newenvironment{slogan}{}{}
\newenvironment{history}{}{}
\usepackage{cleveref}
\crefname{originallemma}{original lemma}{original lemmas}
\Crefname{originallemma}{Original lemma}{Original lemmas}
\makeatletter
\newlabel{lemma-base-change-dimension-fibres}{{lemma-base-change-dimension-fibres}{source}{}{sourcepointer-1}{}}
\newlabel{lemma-base-change-dimension-fibres@cref}{{[originallemma][0][]lemma-base-change-dimension-fibres}{[1][0][]source}{}{}{}}
\makeatother
\crefname{originallemma}{original lemma}{original lemmas}
\Crefname{originallemma}{Original lemma}{Original lemmas}
\makeatletter
\newlabel{lemma-dimension-in-neighbourhood}{{lemma-dimension-in-neighbourhood}{source}{}{sourcepointer-2}{}}
\newlabel{lemma-dimension-in-neighbourhood@cref}{{[originallemma][0][]lemma-dimension-in-neighbourhood}{[1][0][]source}{}{}{}}
\makeatother
\crefname{originallemma}{original lemma}{original lemmas}
\Crefname{originallemma}{Original lemma}{Original lemmas}
\makeatletter
\newlabel{limits-lemma-descend-finite-presentation}{{01ZM}{source}{}{sourcepointer-3}{}}
\newlabel{limits-lemma-descend-finite-presentation@cref}{{[originallemma][0][]01ZM}{[1][0][]source}{}{}{}}
\makeatother
\crefname{originallemma}{original lemma}{original lemmas}
\Crefname{originallemma}{Original lemma}{Original lemmas}
\makeatletter
\newlabel{topology-lemma-characterize-constructible-Noetherian}{{053Z}{source}{}{sourcepointer-4}{}}
\newlabel{topology-lemma-characterize-constructible-Noetherian@cref}{{[originallemma][0][]053Z}{[1][0][]source}{}{}{}}
\makeatother
\title{The level sets of the fibre-dimension function are locally constructible}
\author{The Stacks Project Authors}
\date{stacks-project:more-morphisms.tex:local-snapshot-20261011; GNU FDL 1.2 or later}
\begin{document}
\maketitle
\noindent\textit{Editorial note.} Statement span is lines 8453-8464 of sources/stacks-more-morphisms/pinned-original-source.tex; proof span is lines 8466-8489, the complete author proof (22 body lines).\par\medskip
\noindent\textit{Editorial note.} Navigation stubs: 4 cited-but-not-included labels are supplied as non-load-bearing pointers with their pinned tags where the snapshot tag index has them.\par\medskip
\noindent\textit{Editorial note.} extra\_packages is empty on purpose: the retained preamble already loads xy with options and a second declaration clashes in the assembler wrapper.\par\medskip
\noindent\textit{Editorial note.} Text is retained verbatim, including the author's own typos and placeholders.\par\medskip
\noindent\textit{Allowed standard prerequisites (editorial).}\begin{itemize}
\item Explicit prerequisite boundary. Fibre dimension of a morphism, dimensions of finitely generated algebras, and constructibility; the encoding is retained.
\end{itemize}\medskip
\ifdefined\hypertarget\hypertarget{sourcepointer-1}{}\fi
\noindent\textit{Source reference.} Original lemma lemma-base-change-dimension-fibres: sources/stacks-more-morphisms/pinned-original-source.tex, source lines 8435--8446. Cited by the retained statement or proof but not part of the collected spans; supplied only as a navigation pointer. The argument does not depend on its content.\par\medskip
\ifdefined\hypertarget\hypertarget{sourcepointer-2}{}\fi
\noindent\textit{Source reference.} Original lemma lemma-dimension-in-neighbourhood: sources/stacks-more-morphisms/pinned-original-source.tex, source lines 8412--8417. Cited by the retained statement or proof but not part of the collected spans; supplied only as a navigation pointer. The argument does not depend on its content.\par\medskip
\ifdefined\hypertarget\hypertarget{sourcepointer-3}{}\fi
\noindent\textit{Source reference.} Original lemma 01ZM: sources/stacks-more-morphisms/pinned-original-source.tex, source lines 8489--8489. Cited by the retained statement or proof but not part of the collected spans; supplied only as a navigation pointer. The argument does not depend on its content.\par\medskip
\ifdefined\hypertarget\hypertarget{sourcepointer-4}{}\fi
\noindent\textit{Source reference.} Original lemma 053Z: sources/stacks-more-morphisms/pinned-original-source.tex, source lines 8489--8489. Cited by the retained statement or proof but not part of the collected spans; supplied only as a navigation pointer. The argument does not depend on its content.\par\medskip
\begin{lemma}
\label{lemma-dimension-fibres-constructible}
Let $f : X \to Y$ be a morphism of schemes. Let
$n_{X/Y}$ be the function on $Y$ giving the dimension of fibres of $f$
introduced in
Lemma \ref{lemma-base-change-dimension-fibres}.
Assume $f$ of finite presentation. Then the level sets
$$
E_n = \{y \in Y \mid n_{X/Y}(y) = n\}
$$
of $n_{X/Y}$ are locally constructible in $Y$.
\end{lemma}

\begin{proof}
Fix $n$. Let $y \in Y$. We have to show that there exists an open neighbourhood
$V$ of $y$ in $Y$ such that $E_n \cap V$ is constructible in $V$. Thus we may
assume that $Y$ is affine. Write $Y = \Spec(A)$ and
$A = \colim A_i$ as a directed limit of finite type
$\mathbf{Z}$-algebras. By
Limits, Lemma \ref{limits-lemma-descend-finite-presentation}
we can find an $i$ and a morphism $f_i : X_i \to \Spec(A_i)$ of
finite presentation whose base change to $Y$ recovers $f$. By
Lemma \ref{lemma-base-change-dimension-fibres}
it suffices to prove the lemma for $f_i$. Thus we reduce to
the case where $Y$ is the spectrum of a Noetherian ring.

\medskip\noindent
We will use the criterion of
Topology, Lemma \ref{topology-lemma-characterize-constructible-Noetherian}
to prove that $E_n$ is constructible in case $Y$ is a Noetherian scheme.
To see this let $Z \subset Y$ be an irreducible closed subscheme.
We have to show that $E_n \cap Z$ either contains a nonempty open subset
or is not dense in $Z$. Let $\xi \in Z$ be the generic point. Then
Lemma \ref{lemma-dimension-in-neighbourhood}
shows that $n_{X/Y}$ is constant in a neighbourhood of $\xi$ in $Z$.
This implies what we want.
\end{proof}

\end{document}
